\documentclass[11pt]{article}

\usepackage[margin=1in]{geometry}
\usepackage{amsmath}
\usepackage{tikz}

\title{A Small Reproducible Study of Midpoint Integration}
\author{PI Lab example project}
\date{}

\begin{document}
\maketitle

\begin{abstract}
We approximate $\pi$ with the composite midpoint rule at three resolutions.
The example provides a complete, short paper with a numerical result, a figure,
and a bibliography. It is intended for testing a paper-writing workflow.
\end{abstract}

\section{Method}

The composite midpoint rule is a standard numerical integration method
\cite{atkinson1989}. We apply it to
\begin{equation}
  \pi = 4\int_0^1 \frac{1}{1+x^2}\,dx,
  \qquad
  M_n = \frac{4}{n}\sum_{k=0}^{n-1}
  \frac{1}{1+\left(\frac{k+1/2}{n}\right)^2}.
\end{equation}
The error is $e_n = |M_n-\pi|$. The script
\texttt{analysis/midpoint.py} calculates the values in Table~\ref{tab:results}
with Python's standard library.

\section{Results}

\begin{table}[h]
\centering
\begin{tabular}{rcc}
\hline
$n$ & $M_n$ & $e_n$ \\
\hline
8   & 3.142894730 & $1.3021\times 10^{-3}$ \\
32  & 3.141674034 & $8.1380\times 10^{-5}$ \\
128 & 3.141597740 & $5.0863\times 10^{-6}$ \\
\hline
\end{tabular}
\caption{Midpoint estimates of $\pi$ and absolute errors.}
\label{tab:results}
\end{table}

The error falls by about a factor of 16 when $n$ grows by a factor of four.
Figure~\ref{fig:convergence} shows the same values as correct decimal digits,
defined as $-\log_{10} e_n$.

\begin{figure}[h]
\centering
\begin{tikzpicture}[x=1.6cm,y=0.85cm]
  \draw[->] (0,0) -- (3.2,0) node[right] {$n$};
  \draw[->] (0,0) -- (0,3.35) node[above] {$-\log_{10}e_n$};
  \foreach \y/\label in {0.6/3,1.6/4,2.6/5} {
    \draw (-0.05,\y) -- (0.05,\y);
    \node[left] at (-0.05,\y) {\small \label};
  }
  \foreach \x/\label in {0.5/8,1.5/32,2.5/128} {
    \draw (\x,-0.05) -- (\x,0.05);
    \node[below] at (\x,-0.05) {\small \label};
  }
  \draw[thick] (0.5,0.486) -- (1.5,1.689) -- (2.5,2.894);
  \foreach \x/\y in {0.5/0.486,1.5/1.689,2.5/2.894} {
    \fill (\x,\y) circle (2pt);
  }
\end{tikzpicture}

\caption{Correct decimal digits increase as the midpoint grid is refined.}
\label{fig:convergence}
\end{figure}

\section{Discussion}

These three measurements agree with second-order convergence. They do not
establish the behavior for all $n$, and the comparison uses the floating-point
value of $\pi$ supplied by Python. A larger study could compare a second
integration method and measure the effect of floating-point precision.

\bibliographystyle{plain}
\bibliography{references}

\end{document}
